\chapter{Literature Review}
\label{ch:literature}

This chapter reviews the two architectures the brief asks us to choose
between, and the stream-processing semantics that the choice depends on.
We review them in order to \emph{use} them: every claim made here is
called upon in Chapter~\ref{ch:architecture} or
Chapter~\ref{ch:implementation}.

\section{The Lambda Architecture}
\label{sec:lit-lambda}

Marz and Warren propose Lambda as an answer to a specific
problem: how to compute arbitrary functions over arbitrary data at
scale, when no single system can be simultaneously fast, exact and
fault-tolerant~\cite{marz2015bigdata}. Their construction has three
parts. An \emph{immutable master dataset} records raw facts and is never
updated in place. A \emph{batch layer} recomputes views over the whole
master dataset from scratch. A \emph{speed layer} computes incremental,
approximate views over only the recent data the batch layer has not yet
covered. A serving layer merges the two.

Two properties of that construction matter for this project. First,
\textbf{human fault tolerance}: because the master dataset is immutable
and the batch views are a pure function of it, a bug in view computation
is repaired by fixing the code and recomputing, not by attempting to
correct derived state in place. Helland makes the more general version of
this argument --- that immutability converts a class of distributed-state
problems into append-only bookkeeping~\cite{helland2015immutability}.
Second, the speed layer's error is explicitly \emph{temporary}: it covers
only the window the batch layer has not yet recomputed, and is discarded
when the batch view catches up.

Kleppmann situates the same idea in a wider treatment of derived
data, and makes the point this project leans on most heavily: the
distinction between a system of record and a derived view is what makes
recomputation safe, and losing that distinction is what makes data
pipelines unmaintainable~\cite{kleppmann2017ddia}.

\section{The Kappa Architecture and the Critique of Lambda}
\label{sec:lit-kappa}

Kreps's critique is the standard case against
Lambda and deserves to be taken at full strength~\cite{kreps2014questioning}.
His argument is not that Lambda is incorrect but that it is
\emph{expensive in the wrong currency}. Maintaining two implementations
of the same business logic --- one in a batch framework, one in a
streaming framework --- doubles the code that must be written, tested and
reasoned about, and, worse, produces two implementations that can
silently disagree. Debugging then means debugging the divergence rather
than either system.

His proposed replacement, commonly called Kappa, removes the batch layer
entirely. There is one processing path, a stream processor. Recomputation
is achieved not by a separate batch job but by \emph{replaying the log}
from a chosen offset into a new version of the job, and cutting over when
it catches up. The log, not a file system, is the system of record.

The critique is strong and largely correct as stated. It rests, however,
on an assumption worth naming explicitly, because it is the assumption
this project's use case violates: that \textbf{everything needed to
recompute a result is in the log}. Where that holds, Kappa's replay is a
complete recomputation story and Lambda's batch layer is redundant
machinery. Section~\ref{sec:kappa-rejected} argues that it does not hold
here, and that the failure is structural rather than a matter of tuning.

\section{Event Time, Watermarks and Delivery Semantics}
\label{sec:lit-semantics}

The correctness argument in Chapter~\ref{ch:results} depends on three
results from the stream-processing literature.

\textbf{Event time versus processing time.} Akidau et al.\ formalise the
distinction in the Dataflow model: an event's \emph{event time} is when
it happened; its \emph{processing time} is when the system saw
it~\cite{akidau2015dataflow}. The two diverge arbitrarily under network
delay, retries and partition skew. Any aggregation that groups by
processing time produces results that depend on infrastructure
behaviour rather than on the world, which for a financial figure is
disqualifying.

\textbf{Watermarks are a bet, not a guarantee.} Because a system cannot
know whether more data will arrive for a past window, it must decide when
to stop waiting. A watermark is that decision: a heuristic assertion that
no event older than some bound will now arrive. Dataflow and its
predecessor MillWheel both make the same point --- the watermark trades
completeness against latency, and events later than the watermark are
\emph{dropped}, not deferred~\cite{akidau2015dataflow,akidau2013millwheel}.
This is the mechanism we measure directly in
Section~\ref{sec:consistency}: our watermark is deliberately set tighter
than the injected lateness, so the loss is real and observable rather
than hypothetical.

\textbf{End-to-end exactly-once requires idempotent sinks.} Spark's
Structured Streaming provides exactly-once semantics with respect to its
own state and offsets, through checkpointing and write-ahead
logging~\cite{armbrust2018structured}, building on the micro-batch
fault-tolerance model of discretised
streams~\cite{zaharia2013dstreams} and the lineage-based recovery of
RDDs~\cite{zaharia2012rdd}. That guarantee stops at the sink boundary:
an arbitrary external sink written from \code{foreachBatch} observes
\emph{at-least-once} delivery, because a failure between writing and
committing offsets causes the batch to be re-executed. Kleppmann's
treatment is the clearest statement of the remedy --- make the write
idempotent, typically by keying it, so a repeated delivery is
indistinguishable from a single one~\cite{kleppmann2017ddia}. Every sink
in this project does that, and Section~\ref{sec:impl-sinks} gives the
keys.

\section{The Tooling Landscape}
\label{sec:lit-tooling}

Kafka's original design paper sets out the properties we rely on: a
partitioned, replicated, append-only commit log, with ordering guaranteed
\emph{within} a partition and not across
them~\cite{kreps2011kafka,kafka_docs}. That single sentence is the reason
partitioning is a correctness decision in this project rather than a
throughput one (Section~\ref{sec:impl-partitioning}).

For stream processing, the principal alternative to Spark Structured
Streaming is Apache Flink, whose design treats streaming as primary and
batch as a bounded special case, with a stronger state
model~\cite{carbone2015flink}. We acknowledge in
Section~\ref{sec:tech-stack} that Flink is arguably the better stream
processor considered in isolation, and explain why an engine that also
runs our batch layer wins on a criterion Flink cannot help with.

\section{Positioning of This Work}
\label{sec:lit-gap}

The literature gives a clear decision rule that is, in practice, rarely
applied: choose Lambda when recomputation cannot be expressed as a log
replay, and choose Kappa when it can. Most published comparisons instead
weigh operational simplicity against latency, which is a tie the
fashionable option usually wins.

This project contributes a worked case where the rule bites. The
correcting event --- a partner's revised expense file --- arrives on a
separate channel, in a different format, days after the telemetry it
restates. It is not late data in the watermark sense; it is not in the
log at all. Furthermore, where most treatments assert that the speed
layer is approximate, we \emph{measure} the approximation on a system
configured so that the two layers must disagree, and report the result
including the days on which the disagreement was zero.
